% Generated draft: requires mathematical and source review.
\documentclass[11pt]{article}
\usepackage[margin=25mm]{geometry}
\usepackage{fontspec}
\setmainfont{Latin Modern Roman}
\usepackage{amsmath,amssymb,mathtools,unicode-math}
\setmathfont{Latin Modern Math}
\usepackage{enumitem,needspace,xurl,hyperref}
\setlength{\parindent}{0pt}
\setlength{\parskip}{5pt}
\setlength{\emergencystretch}{3em}
\newcommand{\N}{\mathbb{N}}
\newcommand{\Z}{\mathbb{Z}}
\newcommand{\Q}{\mathbb{Q}}
\newcommand{\R}{\mathbb{R}}
\newcommand{\C}{\mathbb{C}}
\newcommand{\F}{\mathbb{F}}
\newcommand{\A}{\mathbb{A}}
\newcommand{\PP}{\mathbb{P}}
\newcommand{\E}{\mathbb{E}}
\newcommand{\eps}{\varepsilon}
\newcommand{\dd}{\,\mathrm{d}}
\begin{document}
\section*{UnsolvedMath ID: 30006306}
\textbf{Source problem code:} catalogue entry
\subsection{Definitions and mathematical statement}
Let \(M\) be a lattice, meaning a free abelian group of finite rank, and put \(M_{\R}=M\otimes_{\Z}\R\). A lattice polytope is the convex hull of a nonempty finite subset of \(M\). For such a polytope \(P\), write
\[
r=\dim P,\qquad L_P=M\cap\operatorname{span}_{\R}(P-P),\qquad nP=\{nx:x\in P\},\qquad A_n=nP\cap M.
\]
A triangulation of \(nP\) with vertices in \(A_n\) is a finite collection of simplices with vertices in \(A_n\), closed under taking faces, whose union is \(nP\), and such that the intersection of any two simplices is a common face. It is unimodular if, for every maximal simplex \(\operatorname{conv}(v_0,\ldots,v_r)\), the vectors \(v_1-v_0,\ldots,v_r-v_0\) form a \(\Z\)-basis of \(L_P\). For \(r=0\), this uses the empty basis of the zero lattice.

Here a placing triangulation for an ordering \(a_1\prec\cdots\prec a_N\) of \(A_n\) means the ordinary incremental placing construction. Start with the point \(a_1\). At each subsequent point \(a_i\), let \(Q\) be the convex hull of the preceding points. If \(a_i\notin\operatorname{aff}(Q)\), take the cone from \(a_i\) over the current triangulation. If \(a_i\in\operatorname{aff}(Q)\setminus Q\), retain the current triangulation and adjoin cones from \(a_i\) over its triangulated boundary facets strictly visible from \(a_i\), together with all their faces. A boundary facet is strictly visible when \(a_i\) lies strictly on the opposite side of its supporting hyperplane from the relative interior of \(Q\). If \(a_i\in Q\), leave the triangulation unchanged. Thus this convention does not require every ordered point to become a vertex.

\textbf{Proposed statement (catalogue entry).} For every lattice \(M\) and every lattice polytope \(P\subset M_{\R}\), there exist an integer \(n\geq 1\) and a total ordering \(\prec\) of \(A_n\) such that the induced placing triangulation \(T_{\prec}(A_n)\) is unimodular.

\textbf{Unresolved source qualifications.} The supplied record does not define placing triangulation or specify a full-dimensionality hypothesis. The definitions above explicitly use ordinary incremental placing and the intrinsic direction lattice for lower-dimensional polytopes; these conventions require comparison with the cited report. The supplied original excerpt describes a proposed conjectural route that, if true, would imply strong factorization, but does not itself display the polytope statement. Consequently, the exact report formulation and the qualifications on that claimed implication remain to be checked. Strong factorization is background motivation here, not a separately formulated target.
\subsection{Short English statement}
The proposed conjecture asks whether every lattice polytope has a positive integral dilation whose lattice points can be ordered so that incremental placing produces only unimodular maximal simplices. The placing convention and the treatment of lower-dimensional polytopes are explicit above but require source confirmation. This is a statement-only draft requiring human mathematical review.
\subsection{Sources}
\begin{itemize}
\item Supplied local UnsolvedMath registry record: permanent ID \(30006306\), title \emph{Unimodular Placing Triangulations after Lattice Dilation}, source problem code \texttt{catalogue entry}, in the Oberwolfach Reports Open Problems collection. The proposed year is 2025; no individual proposer is identified. The supplied editorial label is \texttt{corrected\_verified}; its accompanying explanation says only formula delimiters, notation, grammar, and typography were repaired while preserving mathematical content and scope. This label is not independent mathematical verification.
\item \emph{Toric Geometry} (2025), Oberwolfach Reports, DOI 10.4171/owr/2025/19, pp.\ 883--940. Supplied report reference and literature-source link: \url{https://doi.org/10.4171/owr/2025/19}. The record attributes the problem to workshop open-problem discussions; an exact passage location and named proposer are not supplied.
\item Adiprasito and Pak, \emph{All triangulations have a common stellar subdivision}. Supplied literature-source link: \url{https://arxiv.org/abs/2404.05930}. The record does not provide a specific result or passage establishing its relationship to the proposed statement.
\item Supplied literature assessment dated exactly \texttt{2026-08-22}: the record labels the problem open, describes it as a proposed sufficient conjecture for strong factorization, says that the report states it without a proof, and describes the underlying smooth strong factorization problem as unresolved. This is supplied dated curation, not an independently verified literature assessment. The supplied background explicitly requests independent verification of time-sensitive status and cited claims.
\end{itemize}
\end{document}
